\section{Related Work}\label{sec:related_work}
\noindent \textbf{KG-based LLM reasoning}.
KGs provide structured knowledge valuable for integration with LLMs~\cite{pan2024unifying}. Early studies~\cite{peters2019knowledge, rogluo2023reasoning, zhang2021poolingformer, li2023chaincok} embed KG knowledge into neural networks during pre-training or fine-tuning, but this can reduce explainability and hinder efficient knowledge updating~\cite{pan2024unifying}.
Recent methods combine KGs with LLMs by converting relevant knowledge into textual prompts, often ignoring structural information~\cite{pan2024unifying, mindmapwen2023}. Advanced works~\cite{tog1.0sun2023think, jiang2023structgpt, tog2.0ma2024think} involve LLMs directly exploring KGs, starting from an initial entity and iteratively retrieving and refining reasoning paths until the LLM decides the augmented knowledge is sufficient.
However, by starting from a single vertex and ignoring the question's position within the KG's structure, these methods overlook multiple topic entities and the explainability provided by multi-entity paths. 
 
% However, these methods often ignore the question's position within the KG's graph structure, starting exploration from a single vertex. This approach overlook the importance of multiple topic entities and the explainability provided by multi-entity paths. fail in long reasoning tasks because the LLM considers only local semantic contexts rather than a global perspective. 
% Moreover, when questions involve multiple entities, they overlook the importance of multiple topic entities and the explainability provided by multi-entity paths. They also typically impose a uniform depth limit, which may not be optimal for all queries.




% \myparagraph{KG-Based LLM Reasoning}KGs offer dynamic, explicit, and structured knowledge representations, making them valuable for integration with LLMs\cite{pan2024unifying}. Early studies \cite{peters2019knowledge, rogluo2023reasoning,zhang2021poolingformer,li2023chaincok} focused on embedding structured knowledge from KGs into neural networks during pre-training or fine-tuning. However, embedding KGs within LLMs can sacrifice explainability in knowledge reasoning and hinder efficient knowledge updating\cite{pan2024unifying}.
% Recent works combine KGs with LLMs by translating relevant structured knowledge into textual prompts, often involving converting the graph into natural language input for the LLM, which ignores the inherent structural information\cite{pan2024unifying,tog1.0sun2023think}. Alternatively, more advanced works \cite{tog1.0sun2023think, jiang2023structgpt,tog2.0ma2024think} involve LLMs in exploring KGs directly, usually identify the initial entity in the question, then to iteratively retrieve and refine the reasoning path until LLM to decide the augmented knowledge is sufficient.
% These existing methods, while effective to some extent, tend to ignore the positional information of the question within the KG's graph structure. They often start exploration from a single vertex, using KG triples as knowledge fed into the LLM to form new reasoning paths. This approach can fail during long reasoning tasks because the LLM considers only local semantic contexts rather than a global perspective. Moreover, when questions involve multiple entities, existing methods overlook the importance of multiple topic entities and the explainability provided by multi-entity paths in question answering. They also typically impose a broad, uniform depth limit for reasoning deep questions, which may not be optimal for all queries.

\myparagraph{Reasoning with LLM prompting}
LLMs have shown significant potential in solving complex tasks through effective prompting strategies. Chain of Thought (CoT) prompting~\cite{wei2022cot} enhances reasoning by following logical steps in few-shot learning. Extensions like Auto-CoT~\cite{zhang2022automatic}, Complex-CoT~\cite{fu2022complexity}, CoT-SC~\cite{wang2022self}, Zero-Shot CoT~\cite{kojima2022large}, ToT~\cite{yao2024ToT}, and GoT~\cite{besta2024graphGoT} build upon this approach. However, these methods often rely solely on knowledge present in training data, limiting their ability to handle knowledge-intensive or deep reasoning tasks.
To solve this, some studies integrate external KBs using plan-and-retrieval methods such as CoK~\cite{li2023chaincok}, RoG~\cite{rogluo2023reasoning}, and ReAct~\cite{yao2022react}, decomposing complex questions into subtasks to reduce hallucinations. 
% However, they may overlook temporal aspects of sub-problems and final answers, potentially leading to locally optimal solutions instead of globally optimal ones.
However, they may focus on the initial steps of sub-problems and overlook further steps of final answers, leading to locally optimal solutions instead of globally optimal ones.
% Our work introduces dynamic multi-hop question reasoning to address deep reasoning challenges. We adaptively determine reasoning depths for different questions, iterating up to a maximum depth, allowing the model to handle varying complexities effectively.
To address these deep reasoning challenges, we introduce dynamic multi-hop question reasoning. By adaptively determining reasoning depths for different questions, we enable the model to handle varying complexities effectively.


% \myparagraph{Existing KG Information Pruning}
% Knowledge graphs store extensive amounts of facts\cite{guo2019learning}, making it impractical to inject all relevant triples into the LLM's context window due to high encoding costs and potential noise\cite{wang2024llm}. Existing work\cite{tog1.0sun2023think, jiang2023structgpt,tog2.0ma2024think} as methioned  typically identifies initial entities and iteratively retrieves and refines reasoning paths until reaching an answer. They often treat the LLM as a function executor, relying on in-context learning or fine-tuning to refine outputs~\cite{dong2022incontext, jiang2024finetune}. However, this heavy reliance on LLMs makes the approach expensive. Some work attempts to address the pruning cost. For instance, KAPING \cite{baek2023knowledge} projects questions and triples into the same semantic space to retrieve relevant knowledge through similarity measures. KG-GPT\cite{kim2023kg} decomposes multi-hop questions into sub-questions, matches relations associated with entities, and selects the top-$k$ relevant relations to form evidence triples. Similarly, KGR\cite{guan2024mitigating} splits retrieved triples into chunks and utilizes LLMs to identify critical triples relevant to the question. However, these pruning methods do not consider the importance of the overall graph structure for pruning and the interrelations among multiple topic entities, leading to suboptimal pruning and reasoning performance.

% Our work addresses these limitations by incorporating graph structural information into the pruning process. We emphasize the importance of multiple topic entities and their structural relationships within the KG, enabling more effective and explainable reasoning paths. By adaptively adjusting reasoning depths and integrating global graph structures, we enhance the LLM's ability to generate faithful and interpretable answers while efficiently narrowing down candidate paths.



\myparagraph{KG information pruning}
Graphs are widely used to model complex relationships among different entities ~\cite{tan2023higher,sima2024deep,lifan2024adarisk,lifan2024hypergraph}.
KGs contain vast amounts of facts~\cite{guo2019learning,DBLP:journals/pvldb/WuSWWZQL24,DBLP:journals/pvldb/WangWWZQZL24}, making it impractical to involve all relevant triples in the context of the LLM due to high costs and potential noise~\cite{wang2024llm}. Existing methods~\cite{tog1.0sun2023think, jiang2023structgpt, tog2.0ma2024think} typically identify initial entities and iteratively retrieve reasoning paths until an answer is reached, often treating the LLM as a function executor and relying on in-context learning or fine-tuning, which is expensive.
Some works attempt to reduce pruning costs. KAPING~\cite{baek2023knowledge} projects questions and triples into the same semantic space to retrieve relevant knowledge via similarity measures. KG-GPT~\cite{kim2023kg} decomposes complex questions, matches, and selects the relevant relations with sub-questions to form evidence triples. 
% Similarly, KGR~\cite{guan2024mitigating} splits retrieved triples into chunks and uses LLMs to identify critical ones. 
% retrofitting the initial draft responses of LLMs based on the
% factual knowledge stored in KGs.
% However, these methods often neglect the overall graph structure and interrelations among multiple topic entities, leading to suboptimal pruning and reasoning performance.
However, these methods often overlook the overall graph structure and the interrelations among multiple topic entities, leading to suboptimal pruning and reasoning performance.
% Our work addresses these limitations by incorporating graph structural information into the pruning process. We emphasize the importance of multiple topic entities and their structural relationships within the KG, enabling more effective and explainable reasoning paths. By adaptively adjusting reasoning depths and integrating global graph structures, we enhance the LLM's ability to generate faithful and interpretable answers while efficiently narrowing down candidate paths.
% We address these limitations by incorporating graph structural information into the pruning process. By emphasizing multiple topic entities and their relationships within the KG, we enable more effective and explainable reasoning paths. Adaptively adjusting reasoning depths and integrating global graph structures enhance the LLM's ability to generate accurate and interpretable answers while efficiently narrowing down candidate paths.